% =====================================================================
% tesi/capitoli/01-bit-e-numeri.tex
% =====================================================================
% copre: docs/02-numeri-e-bit.md
% ---------------------------------------------------------------------

\chapter{La rappresentazione dei numeri, e gli errori che non fanno rumore}
\label{cap:bit}

Ogni errore descritto in questo capitolo condivide una proprietà, ed è quella proprietà, più che la difficoltà tecnica, a giustificare il fatto che il capitolo apra la trattazione: nessuno di questi errori si manifesta come un guasto. Un indirizzo sbagliato produce di norma un arresto del programma oppure un valore fuori scala, e in entrambi i casi chi osserva se ne accorge. Un byte letto nell'ordine opposto a quello con cui è stato scritto, o un gruppo di quattro bit estratto dalla metà sbagliata di un byte, producono invece un numero perfettamente plausibile: un livello resta un livello, un identificativo resta un identificativo, una statistica resta nell'intervallo delle statistiche. Il difetto attraversa quindi indenne ogni controllo condotto per ispezione, e si rivela soltanto molto più tardi, quando il dato viene confrontato con la sua origine oppure quando il gioco che lo consuma si comporta in un modo che nessuno ha scelto.

Da questa asimmetria fra la gravità del difetto e la sua visibilità discende il criterio metodologico che governa l'intero lavoro, e che va enunciato subito perché regola tutto ciò che segue: nessuna costante di formato viene trascritta a mano quando esiste una fonte da cui possa essere ricavata da un programma. La giustificazione empirica di quel criterio, cioè la misura di quanto le fonti secondarie sbaglino su questi valori, è documentata nel capitolo~\ref{cap:testo}.

\section{L'ordine dei byte}

Un numero che non entra in un solo byte occupa due o quattro byte consecutivi in memoria, e la sua disposizione richiede una scelta che nessun principio impone: collocare per primo il byte più significativo, oppure quello meno significativo. La prima convenzione si chiama \term{big-endian}, la seconda \term{little-endian}. La proprietà che rende la scelta insidiosa è che le due convenzioni non sono distinguibili osservando i byte: la stessa sequenza è valida in entrambe le letture e restituisce due numeri diversi, entrambi legittimi.

Le due famiglie di sistemi trattate qui stanno sui lati opposti di questa convenzione. Le generazioni 1 e 2 girano su un processore Sharp LR35902, derivato dallo Z80, e conservano i numeri a più byte in ordine big-endian: la coppia \hx{01} \hx{F4} vale cinquecento. La generazione 3 gira su un ARM7TDMI e adotta l'ordine little-endian: la medesima coppia di byte vale sessantaduemilaquattrocentosessantacinque. Ne segue che un lettore scritto per la generazione 1 e riusato sulla generazione 3 senza intervenire sull'ordine non produce un errore diagnosticabile, ma un insieme coerente di numeri sbagliati.

Esiste in questo quadro un caso limite che va registrato esplicitamente, perché a una prima lettura appare come una contraddizione della documentazione anziché come una proprietà del codice originale: i checksum della generazione 2 sono valori a sedici bit memorizzati in ordine little-endian, all'interno di una generazione che per ogni altro campo è big-endian. La verifica sul disassemblato \cite{pokecrystal} conferma che si tratta di una scelta del codice originale e non di un errore di trascrizione delle fonti, e la conseguenza pratica è che chi scrive un writer deve trattare quel campo diversamente da tutti gli altri campi della stessa struttura, cioè deve sospendere per un solo campo la regola che ha appena applicato a tutti i precedenti.

\section{Il nibble, ovvero la metà di un byte}

Un byte contiene otto bit, e le architetture di quell'epoca non tolleravano che un solo bit rimanesse inutilizzato. Da questo vincolo economico derivano due tecniche di impacchettamento che ricorrono in ogni struttura descritta in questo lavoro.

La prima è il \term{nibble}, cioè un gruppo di quattro bit, capace di rappresentare un valore fra zero e quindici. Le generazioni 1 e 2 memorizzano in due soli byte i quattro valori individuali di un Pokémon proprio in questa forma: il primo byte porta l'Attacco nel nibble alto e la Difesa nel nibble basso, il secondo la Velocità nel nibble alto e la statistica Speciale nel nibble basso \cite{pokered}. L'estrazione del nibble alto consiste in uno scorrimento a destra di quattro posizioni, quella del nibble basso in una mascheratura con \hx{0F}, e lo scambio involontario delle due operazioni costituisce il caso esemplare dell'errore silenzioso: entrambi i valori estratti restano nell'intervallo ammesso, e nulla nel dato segnala l'inversione.

La seconda tecnica è il \term{campo di bit}, cioè un valore numerico che occupa un intervallo di bit non allineato al confine del byte. La sua realizzazione più densa in questo progetto è la parola da trentadue bit della generazione 3 che ospita sei valori individuali da cinque bit ciascuno, il contrassegno di uovo e lo slot di abilità: otto informazioni distinte in quattro byte \cite{pokeemerald}. La lettura di un campo di bit richiede uno scorrimento a destra pari al numero di bit che lo precedono e una mascheratura con tanti bit a uno quanti il campo ne occupa; la sua scrittura richiede l'azzeramento della sola finestra interessata e l'inserimento del valore al suo interno, senza alterare i campi adiacenti, il che significa che scrivere un campo di bit è un'operazione di lettura, modifica e riscrittura e non una semplice assegnazione.

Un dettaglio di questa meccanica produce errori con particolare frequenza, e conviene enunciarlo come regola: quando i bit di un campo attraversano il confine fra due byte, l'ordine dei byte torna a contare. La conseguenza operativa è che l'ordine corretto delle operazioni consiste nel leggere l'intera parola rispettando l'endianness dichiarata dall'architettura, e soltanto dopo operare sui bit del numero così ottenuto; l'ordine opposto, cioè operare byte per byte e comporre alla fine, funziona nei casi allineati e fallisce esattamente nei casi in cui il campo attraversa il confine.

\section{Il byte che contiene due informazioni distinte}

Una variante del campo di bit merita una trattazione separata, perché in essa due informazioni concettualmente indipendenti condividono un byte per sola economia di memoria, e perché la loro separazione è ciò che una conversione fra generazioni deve compiere. Nelle generazioni 1 e 2 il byte che descrive i punti potenza di una mossa impiega i sei bit bassi per i punti correnti e i due bit alti per il numero di incrementi permanenti applicati a quella mossa \cite{pokered}. In generazione 3 le medesime due informazioni risiedono in strutture diverse: i punti potenza nella sottostruttura delle mosse, i bonus nella sottostruttura della crescita \cite{pokeemerald}.

La conversione deve dunque smontare un byte e distribuirne le due metà su due destinazioni distinte, e questo è il primo esempio, fra i molti che il capitolo~\ref{cap:conversione} raccoglie, di un'operazione che non è una copia di byte ma una trasformazione del modello dei dati. Chiamarla copia è già l'errore.

\section{Perché una costante si genera invece di trascriverla}

La conclusione operativa di questo capitolo è un metodo e non una tabella. Ogni volta che un valore costante può essere ricavato da una fonte autorevole, va ricavato da quella fonte per mezzo di un programma, e non ricopiato a mano nel codice. La formulazione appare pedante fino al momento in cui si misura quanto costi la sua violazione, e la misura in questo progetto esiste: la tabella di codifica dei caratteri era errata in due punti in tutte le fonti secondarie consultate, e lo scostamento è emerso soltanto dal confronto con i file di definizione dei caratteri del disassemblato. La risposta a quell'episodio non è stata la correzione manuale della tabella, che avrebbe potuto essere sbagliata di nuovo a mano, ma la sua generazione a partire dalla fonte, per mezzo di uno strumento che rifiuta di emettere il file quando alcuni valori di controllo non corrispondono. Lo strumento e la sua logica di rifiuto sono descritti nel capitolo~\ref{cap:testo}.
